\section{第四条线：为什么速率相近，闭合却仍然是核心问题}\label{story:closure}
GEM式宏观一致性留下一个更严格的问题：模型能否同时预测密度、压力、电流层、热流和不稳定性？Shay的跨模型密度空穴已经显示，相同输出通量可由不同宽度和峰速实现。本轮F611把这一点推进到电子各向异性：流体电子状态方程改变延长电流层和面外磁场，后期总率却仍相近。

由此产生两条互补路线。一条在明确俘获与磁化条件下构造各向异性状态方程；另一条演化完整压力张量，并用动理学热流检验局地、非局地或梯度闭合。Landau与分群电子模型又分别补入沿场相混合和多来源分布信息。以下按“增加什么信息—改善哪个可测量量—还错在哪里”比较，不把增加矩数或一个漂亮的岛图当成闭合已成功。

\subsection{密度空穴是额外检验：总率相近仍可能有显著结构误差}
Shay等2001年的跨模型研究不仅比较总率，还明确报告分离线密度耗尽：相同上游密度0.2下，Hall流体最低约0.03、混合模型约0.1、PIC约0.15。离子动理学喷流约宽两倍而峰速约减半，总输出通量因而仍接近。原文将空穴联系到Hall电流造成的磁场堆积和排斥离子的电场；有限轨道效应减弱空穴。\citep{Shay2001}

因此，“Hall模型算对了重联率”不能单独用来认证卫星沿轨道密度峰谷。该流体计算还为防止负密度加入二阶扩散，电子惯性用常$d_e$反演，使电子层比当地低密度PIC约窄两倍。其无质量电子计算靠六阶耗散断线，不是理想Hall模型自行改变磁拓扑。

该文关于小均匀电阻不能平衡电场的论证，使用了电子层宽长比近固定的几何关系。后续电阻Hall层理论放松几何约束，得到不同分支。阅读这些结论时，应先对照层长怎样决定，再判断是否互相矛盾。\citep{Shay2001,Malyshkin2008,Simakov2009}

\subsection{各向异性状态方程：长电子层出现，后期总率仍相近}\label{sec:synthesis-anisotropy}
\citet{F611}在中等导引场的force-free和Harris片中比较全PIC、各向异性电子混合模型和等温电子混合模型。最直接的新现象是：前两者出现跨数个$d_i$的电子电流层，等温例没有同样长的层；Harris片X线附近的面外磁场还增强到背景的约1.5--2倍。图6的沿片剖面为双峰，不能把所有增强都画成X点正中心的单峰。

这补足了前面Sullivan层长研究与压力闭合研究之间的联系：层长不仅会随耗散改变，也可能由电子各向异性支撑。force-free例的$F_e=\mu_0(p_{e\parallel}-p_{e\perp})/B^2$峰值约0.7；Harris长层处约大于0.5。但图3、7的后期率在三种模型中仍接近，说明\textbf{复现总率不等于复现长层及Hall场}。早期时钟并不一致，Harris比较还用了不同种子并作时间平移。

这一结果与\citet{Cassak2015}不能合并成“任意电子各向异性均能独自产生快重联”。两者的状态方程不同；F611明确指出，所用俘获插值闭合在该参数区主要改变结构，并不自行提供无Hall快支。其Harris实现还要沿二维磁通标签追踪上游值，并修改闭合的参考密度以避免中心非物理各向异性。适用于无导场和复杂三维几何的可迁移闭合仍未建立。第一章第\ref{sec:synthesis-closure-model}节给出状态方程、极限和曲率电流的逐步推导。

\subsection{梯度闭合的推进：从直接热流检验到三维不稳定性}
\citet{L256}先从Vlasov分布函数计算热流矩，检查局地各向同性松弛闭合是否给出正确的散度形状、符号和幅度，再提出对压力偏离各向同性部分作Laplacian的闭合。这一步比仅比较重联率更有约束力：局地松弛会在实际热流很弱的上游区域产生错误的热流散度，改变 $k_0$ 不能同时修正所有分量的结构。

其梯度闭合在磁岛尺度 $\lambda/d_i=5,10,15,25$ 的算例中产生磁岛反弹，平均率约随 $\lambda^{-0.73}$ 下降；两套分辨率得到接近的平均与峰值率。但是GEM、大Harris片与岛合并分别用了 $k_{0s}d_s=3,1/3,1/2$。作者明确按动理学结果选取参数，故这些结果证明了模型的拟合与迁移潜力，还没有证明无需校准即可预测所有电流片。

\citet{E552}进一步使用统一的 $\chi=3,k_{0s}d_s=1$ 检验重联、岛合并和低杂漂移不稳定性（LHDI）。质量比100的特定静电LHDI算例中，Vlasov、梯度十矩、五矩和局地松弛十矩的增长率分别约为 $2.2,2.4,3.2,0.67\Omega_i$。这显示增加矩数本身不是充分条件：\textbf{五矩可以把不稳定性发展得过快，某些十矩闭合又会抑制过强。}最快增长模式会随模型和质量比变化，只检查预先指定的单一波数可能得出不同判断。

该文电磁LHDI支的Vlasov／流体增长率约为 $0.41/0.21\Omega_i$，都出现了超过初始重联场量级的磁扰动与强电流片弯曲，其高饱和原因仍不清楚。三维重联较二维更早启动，后期电流片加宽，出流区湍动强于电子扩散区；后期另一个不稳定性仅被称为“类似电流片剪切不稳定性”，尚未完成模态确诊。

一个尤其值得研究的新情况是\textbf{噪声决定不同过程的先后顺序}：密度扰动幅度由 $10^{-10}$ 增至 $10^{-2}$，独立测试的LHDI启动约从 $30\Omega_i^{-1}$ 提前到 $4\Omega_i^{-1}$。若不稳定性先长成，它能建立多个重联点；若种子重联先启动，又会改变密度梯度并减弱LHDI。因此，比较流体与PIC时，粒子噪声、初始种子与闭合导致的启动延迟都必须记账，不能把所有湍动差异都归因于“二维还是三维”。

\subsection{十矩三维闭合：恢复不稳定性，同时产生新的误差}
\citet{Wang2015}的十矩与PIC比较不能概括成时间历程一致。两者在不同时间达到相同重联通量时，近X点一些结构相近；但速率随时间的曲线不同，某压力非对角分量峰值约相差四倍，下游压力结构也不一致。

\citet{Closure2026}在基于MMS非对称事件的三维设置中比较局地与梯度闭合。梯度闭合恢复了局地闭合所压制的低混合漂移和弯曲模，并产生次级磁岛、磁通绳和混合。计算采用质量比100、$1152\times576\times288$ 网格，梯度闭合开销约为局地方案三倍。

同时，梯度模型仍产生范围过大的各向异性与非旋转对称压力；部分峰值出现在与MMS／PIC不同的一侧；弯曲模在Vlasov中数个回旋周期即可发展，而该流体算例需数十周期。文章未提供网格与质量比联合扫描。它说明闭合会改变哪些模式能够出现，并未证明十矩已与动理学等价。

\subsection{加入Landau闭合：岛面积相近，增长时钟与温度结构已改变}
\citet{L244}构造含电子平行／垂直温度和离子垂直温度的六场模型，实际重联模拟则使用冷离子子系统。保留非局地电子热流后，线性增长下降；某些较高斜向模数的模式由不稳定变为稳定。闭合在谐振面局地为零，不妨碍它通过面外响应改变整个本征模。

非线性图像若用各自的$\gamma t$作横轴，岛形和最终面积可以相近，但实际时间已经不同。不能把这种重标度后的重合写成“Landau阻尼不改变重联速度”。温度、密度和电流的细结构也变平滑，所示电子平行温度峰值约降为无热流情况的一半。新发现是\textbf{宏观岛面积对闭合相对稳健，细结构和热响应却不稳健}。

其热流自由能损失可作为向未解析速度空间转移的代理，但模型没有复现所引用动理学研究中磁能释放与电子加热之间的明显时间延迟；作者提出闭合可能耗散过强。另一个量化限制是，约$0.002H_0$的数值能量误差已占本次释放磁能约20\%。故能量分配的精确比例没有像磁岛形态那样受到充分验证。

这些算例只含单一螺旋性，可以用$y^*=y+(k_z/k_y)z$约成二维；它们验证斜向结构与沿场闭合，却不能等同于多个斜向模相互作用的完整三维湍流。原文图4--7和能量段应一起读，避免只选最相似的岛图而忽略时间、闭合耗散与能量误差。

\subsection{来源分群的新闭合方向：改善形状与解决热流并非同一结论}
2026年McGrae-Menge等预印本把片两侧电子分别演化为十矩流体。每群自身三阶矩较小时，群间相对漂移和温度差仍能产生很大的\emph{总}热流；第一章式\eqref{eq:mixture_heat}给出精确分解。与单电子十矩相比，两群模型改善了电子层密度空洞和压力堆积，重建的热流符号与空间结构更接近PIC，但幅度仍偏高。\citep{F635}

其证据边界尤其重要：离子固定；参数精度图是对各PIC算例在时间内选取\emph{最佳}重建指标，不是两流体模拟的全时段平均精度；不同图使用独立色标。高$\beta$时人工标签梯度可引发错误膨胀，晚期同群内部又出现多束。论文提出的动态分裂、三维、湍流及离子耦合尚未验证。所谓156倍压缩是表示自由度估算，不是测得的运行加速。这是一条具体且值得检验的新方向，尚不能写成一般热流闭合问题已经解决。

\paragraph{这一线索走到了哪里。}这条线的进展是逐步增加可独立检验的量：先率与岛形，再压力散度与层长，随后是热流、不稳定模、启动延迟和三维混合。未解决的是同一组闭合参数能否跨Harris片、岛合并、非对称和三维环境迁移。模型应同时通过这些约束；改善某一个结构而破坏另一个结构，需要报告为取舍，不能只展示改善项。
